\documentclass[preview]{standalone}

\usepackage{amsmath}
\usepackage{amssymb}
\usepackage{stellar}
\usepackage{definitions}

\begin{document}

\id{psicologia-stress-coping-resilienza}
\genpage

\section{Stress, coping e resilienza}

\begin{snippetdefinition}{stress-definition}{Stress}
    Lo \emph{stress} è un elevato tasso di attivazione fisica e psichica della
    persona in risposta a eventi interni ed esterni che la coinvolgono.
\end{snippetdefinition}

\begin{snippet}{stress-organismo}
    Lo stress coinvolge diversi sistemi e apparati dell'organismo, come sistema
    respiratorio, cardiaco e urogenitale. Non deve essere inteso
    necessariamente come esperienza negativa: può sostenere risposte costruttive
    oppure diventare destabilizzante.
\end{snippet}

\begin{snippetdefinition}{sindrome-generale-adattamento-definition}{Sindrome generale di adattamento}
    La \emph{sindrome generale di adattamento}, o \emph{SGA}, è il modello che
    descrive lo sviluppo delle risposte dell'individuo di fronte a eventi
    stressanti.
\end{snippetdefinition}

\begin{snippet}{fasi-sindrome-generale-adattamento}
    La sindrome generale di adattamento comprende tre fasi:
    \begin{enumerate}
        \item \emph{allarme}, in cui l'individuo diventa consapevole della
        novità della situazione;
        \item \emph{resistenza}, in cui l'evento stressante perdura e
        l'organismo cerca di adattarsi;
        \item \emph{esaurimento}, in cui si verifica una riduzione della
        resistenza.
    \end{enumerate}
\end{snippet}

\begin{snippetdefinition}{eustress-definition}{Eustress}
    L'\emph{eustress} è un processo di attivazione in cui la persona, di fronte
    a un evento che richiede l'utilizzo di energie, risponde in modo costruttivo.
\end{snippetdefinition}

\begin{snippetdefinition}{distress-definition}{Distress}
    Il \emph{distress} è un processo conseguente a un evento stressante percepito
    come minaccioso e destabilizzante, rispetto al quale l'individuo si sente
    privo di mezzi di difesa.
\end{snippetdefinition}

\begin{snippetdefinition}{coping-stress-definition}{Coping}
    Il \emph{coping} è l'insieme dei meccanismi di fronteggiamento di eventi
    interni ed esterni alla persona, percepiti come minacciosi.
\end{snippetdefinition}

\begin{snippet}{strategie-coping}
    Le strategie di coping possono essere \emph{emotion-focused}, quando mirano
    a modificare le risposte emotive interne, oppure \emph{problem-focused},
    quando consistono in tentativi di risoluzione del problema.
\end{snippet}

\includesnpt[width=58\%,src=/snippet/static-psychology/stress-coping-resilience.jpg]{centered-img}

\begin{snippetdefinition}{resilienza-definition}{Resilienza}
    La \emph{resilienza} è uno stile di coping che consiste nella capacità di
    gestire positivamente eventi dolorosi che si protraggono a lungo.
\end{snippetdefinition}

\begin{snippet}{stress-coping-resilienza}
    Stress, coping e resilienza descrivono tre livelli della risposta alle
    difficoltà: l'attivazione psicofisica, le strategie di fronteggiamento e la
    capacità di mantenere o recuperare un funzionamento positivo nonostante
    eventi prolungati o dolorosi.
\end{snippet}

\end{document}
